\section{A quantitative strain gap and the failure of fixed-control approximation on the supplied fan}
\label{sec:r13-killing-stress-l1-gap}

Keep the five planes, trace data, fields \(b_i\), strains \(Z_i\),
and edge stress \(s\) of PK.  Equip \(N=\mathbf R^4\) and its
subspaces with their Euclidean norms.  Write
\[
 B_R^i=P_i\cap B_R(0),\qquad
 D_R^{ij}=S_{ij}\cap B_R(0),\qquad
             \mathfrak s(c)=\sum_{i<j}s_{ij}c_{ij}.
 \tag{KG1}
\]
All tensor-measure total variations below use the Euclidean
Frobenius norm.  Changing to the fixed lattice calibrations
changes only the constants.

The two external analytic inputs are the classical
Korn--Poincare estimate in \(\mathrm{BD}\) and the bounded
\(\mathrm{BD}\) trace map on a bounded Lipschitz domain.
Here \(Eu=\operatorname{sym}Du\) denotes the full distributional
symmetric derivative.  In the form needed on a three-ball,
there is an affine Killing field \(K\) such that
\[
 \|u-K\|_{L^1(B_R)}\leq C R\,|Eu|(B_R).
 \tag{KG2}
\]
The trace estimate on a half-ball is
\[
 \|\gamma u\|_{L^1(\partial B_R^+)}
       \leq C\bigl(R^{-1}\|u\|_{L^1(B_R^+)}
                              +|Eu|(B_R^+)\bigr).
 \tag{KG3}
\]
We use KG2 only for continuous \(\mathrm{BD}\) fields on a
closed ball; hence they are also in \(L^2\).  KG2 follows by
Holder and scaling from the \(L^{3/2}\) estimate stated as
Theorem A.1 in Friedrich's paper cited below.  KG3 follows
by scaling from Babadjian's Theorem 3.2; for a continuous
field its trace is its usual boundary restriction.  No
\(L^1\) estimate for the full gradient is used.

\subsection{The finite Killing stress is stable in the strain-measure norm}

\begin{theorem}[Quantitative obstruction]
\label{thm:r13-strain-measure-gap}
There is a constant \(C>0\), depending only on the fixed
five-plane configuration, such that the following holds.
Let \(g_i\) be continuous covector fields on
\(\overline{B_R^i}\), belonging to \(\mathrm{BD}(B_R^i)\).
Assume that their tangential traces agree on every pair disk:
\[
                  g_i(x)|_{S_{ij}}=g_j(x)|_{S_{ij}}
                           \quad(x\in D_R^{ij}).
 \tag{KG4}
\]
Then
\[
 R^3|\mathfrak s(c)|
       \leq C\sum_{i=0}^4
             |Z_i\,\mathcal L^3-Eg_i|(B_R^i).
 \tag{KG5}
\]
In particular, when \(c_{01}=1\) and the other edge data
are zero, the sum on the right has a strictly positive lower
bound proportional to \(R^3\).  The conclusion applies to
the restrictions of any single smooth ambient covector field.
\end{theorem}

\begin{proof}
First take \(R=1\).  Put
\(u_i=b_i-g_i\) and
\(\eta_i=|Eu_i|(B_1^i)\).  The full identity
\(Eb_i=Z_i\mathcal L^3\) gives exactly the error measure
appearing in KG5.  By KG2 choose
\[
 K_i(x)=k_i+\Omega_i(x,\cdot),\qquad
              \Omega_i\in\bigwedge^2P_i^*,
 \qquad r_i=u_i-K_i,
\]
with \(\|r_i\|_{L^1(B_1^i)}\leq C\eta_i\).
Its symmetric derivative is still \(Eu_i\).
Apply KG3 to either half-ball cut out by \(S_{ij}\).
Restriction can only decrease the \(L^1\) norm and the
variation of the derivative measure, so
\[
       \|r_i|_{D_1^{ij}}\|_{L^1(D_1^{ij})}
                           \leq C_{ij}\eta_i.
 \tag{KG6}
\]
Continuity identifies this trace from either side.  The
finitely many constants can be replaced by their maximum.

Let \(\widehat\theta_{ij}\) be the alternating form on
\(S_{ij}\) with
\(\widehat\theta_{ij}(a_{ij},b_{ij})=1\).  Denote the
coefficient vectors of \(\Omega_i\) by \(\omega_i\), as in
PK2--PK5.  On \(D_1^{ij}\), the common \(g\) trace cancels.
Using PK9 gives the exact identity of tangential covectors
\[
 (r_i-r_j)|_{S_{ij}}
 =\bigl(c_{ij}-(D\omega)_{ij}\bigr)
                 \widehat\theta_{ij}(x,\cdot)
                    -(k_i-k_j)|_{S_{ij}}.
 \tag{KG7}
\]
The disk is symmetric under \(x\mapsto-x\).  For any
linear map \(L\) and constant vector \(a\), the triangle
inequality gives
\[
       \int_{D_1^{ij}}|Lx|\,dx
                     \leq\int_{D_1^{ij}}|Lx-a|\,dx.
 \tag{KG8}
\]
Indeed integrate
\(2|Lx|\leq|Lx-a|+|-Lx-a|\).
The constant
\(\kappa_{ij}=\int_{D_1^{ij}}
                  |\widehat\theta_{ij}(x,\cdot)|\,dx\)
is strictly positive.  KG6--KG8 therefore imply
\[
       |c_{ij}-(D\omega)_{ij}|
                          \leq C'(\eta_i+\eta_j).
 \tag{KG9}
\]
Pair this inequality with the finite stress PK6.  Since
\(s^TD=0\), it gives
\( |\mathfrak s(c)|\leq C\sum_i\eta_i\).
There is no bound or convergence assumption on the translations
\(k_i\) or rotations \(\Omega_i\); the symmetric-disk
argument removes the translations and the stress removes
all rotations.

For general \(R\), replace \(g_i(x)\) on the unit ball by
\(R^{-1}g_i(Rx)\).  The target \(b_i\) is homogeneous of
degree one, and the variation of its strain error rescales
by \(R^{-3}\).  The unit-ball result yields KG5.
\end{proof}

\subsection{Smooth full-exact forms are included}

Let \(a=d'\beta+d''\zeta\) be a smooth real \((1,1)\)-form
on an ambient neighborhood of the closed ball.  Write
\(\beta=\sum p_j\,d''x_j\) and
\(\zeta=\sum q_j\,d'x_j\).  In the usual superform signs,
\[
 a_{ij}=\partial_i p_j-\partial_j q_i,
 \qquad
 \operatorname{sym}a=e(B),
       \qquad B=\sum_j(p_j-q_j)\,dx_j.
 \tag{KG10}
\]
Restriction to a linear plane commutes with this identity.
Consequently KG5 gives
\[
 \sum_i\int_{B_R^i}
       |Z_i-\iota_i^*\operatorname{sym}a|\,d\mathcal L^3
                    \geq C^{-1}R^3|\mathfrak s(c)|.
 \tag{KG11}
\]
This lower bound holds for all such \(a\), without bounds
on their coefficients or derivatives.

\subsection{A fixed finite positive current cannot dominate vanishing errors}

We make the current-order conclusion precise.  Work in an
ambient open neighborhood \(O\) of \(\overline{B_R(0)}\).
For a smooth full-exact \(a\), form the fixed-fan current
\[
 \mathcal E_a=\sum_i\iota_{i*}
       [Z_i-\iota_i^*\operatorname{sym}a].
 \tag{KG12}
\]
Let \(\mathcal Q\) be a fixed real current of bidegree
\((2,2)\) which is nonnegative on compactly supported
strongly positive smooth real \((2,2)\) tests.  For each
fixed constant strongly positive test form, its evaluation
on \(\mathcal Q\) is thus a locally finite positive Radon
measure.  Say that
\(-\epsilon\mathcal Q\leq\mathcal E_a\leq
\epsilon\mathcal Q\) when both differences are
nonnegative on all these tests.  This definition is sufficient
for the result; no identification of the strong, middle and
weak positivity cones is made.

\begin{corollary}[No vanishing fixed-control error on PK10]
\label{cor:r13-no-fixed-control-on-supplied-five-plane-fan}
If \(\mathfrak s(c)\ne0\), there is no sequence of smooth
full-exact \((1,1)\)-forms \(a_j\) and numbers
\(\epsilon_j\downarrow0\) satisfying
\[
             -\epsilon_j\mathcal Q\leq\mathcal E_{a_j}
                                      \leq\epsilon_j\mathcal Q
 \tag{KG13}
\]
on \(O\), for one fixed \(\mathcal Q\) as above.
\end{corollary}

\begin{proof}
For each \(P_i\), choose an orthonormal basis of its
three-dimensional tangent space.  The six scalar quadratic
probes in the directions \(e_a\) and \(e_a+e_b\),
\(1\leq a<b\leq3\), determine every symmetric matrix
entry by polarization.  For each direction choose a constant
decomposable positive ambient \((2,2)\)-form whose restriction
to \(P_i\) pairs with that quadratic probe, including its
fixed positive normalization.  Such a form is obtained from
two covectors spanning the perpendicular two-plane and their
double-primed copies.  Write these thirty forms as
\(\Psi_{i\nu}\).

For each \(\Psi_{i\nu}\), KG13 dominates the corresponding
scalar signed measure by
\(\epsilon_j\,\mathcal Q(\Psi_{i\nu})\).
Let
\(A_i=P_i\setminus\bigcup_{k\ne i}P_k\).
The five intrinsic three-dimensional error measures are
mutually singular: their intersections have three-dimensional
measure zero.  Restricting the scalar measure inequality to
\(A_i\) therefore isolates exactly the desired scalar
probe of the \(i\)-th strain error.  No other source plane
contributes on this Borel set.

Choose a nonnegative compact cutoff \(\chi\) in \(O\)
which is one on \(\overline{B_R(0)}\).  Polarization and
finite-dimensional norm comparison give a constant \(C_0\)
such that
\[
 \sum_i\int_{B_R^i}
       |Z_i-\iota_i^*\operatorname{sym}a_j|\,d\mathcal L^3
 \leq C_0\epsilon_j
       \sum_{i,\nu}\langle\mathcal Q,\chi\Psi_{i\nu}\rangle.
 \tag{KG14}
\]
Every term in the finite sum is finite and nonnegative,
and the sum is independent of \(j\).  Its right side tends
to zero, contrary to KG11.
\end{proof}

\paragraph{What this rules out.}
The same fixed supplied fan has all the local primitives and
fixed-fan \(d'd''\)-closure proved in PK, yet even arbitrarily
large smooth ambient full-exact coefficients cannot approach
its strain in the above fixed-positive-control order.  This
is a quantitative obstruction on retained real source planes.
No original analytic zero realizing PK10 has been constructed,
and KG does not refute the original analytic all-refinement
criterion.  Genuine actual unit refinements can change the
incident spaces and the associated stress; that operation is
not represented by a real subdivision of this fixed fan.

\paragraph{Primary analytic sources.}
M. Friedrich, \emph{A Korn--Poincare-type inequality for special
functions of bounded deformation}, arXiv:1503.06755v3,
Theorem A.1, states the classical \(\mathrm{BD}\)
Korn--Poincare estimate used in KG2 in arbitrary dimension.
J.-F. Babadjian, \emph{Traces of functions of bounded deformation},
arXiv:1308.5497v2, Theorem 3.2, proves the bounded trace map
on Lipschitz domains and its agreement with continuous
boundary restrictions.  Only these two analytic estimates
are imported; KG5--KG14 are the finite-configuration and
current-order deductions given here.
